\subsection{Proof of Proposition~\ref{prop:unique} (uniqueness of the minimal parent
set) and of Proposition~\ref{prop:welldef}}
\label{proof:P2}
\textbf{Setting.} All variables are finite-valued functions of $X\sim P$.
$U=\{U_1,\dots,U_n\}$ is the candidate parent set; for $A\subseteq U$ write
$A$ also for the tuple of its members. \emph{Positivity}: for every
combination $u=(u_1,\dots,u_n)$ in $\prod_j\mathcal{V}_j$, where
$\mathcal{V}_j$ is the set of values $U_j$ takes with positive probability,
$P(U=u)>0$. Positivity of $U$ implies positivity of every sub-collection,
since a marginal of a positive law on a product set is positive on the
product of the marginal value sets. A set $S\subseteq U$ is a \emph{blanket}
for $C$ if $C\perp U\setminus S\mid S$.

\textbf{Lemma (intersection property with a deterministic target).} Let
$Y,W,Z$ be disjoint sub-collections of $U$ and assume positivity. If
$C\perp Y\mid(Z,W)$ and $C\perp W\mid(Z,Y)$, then $C\perp(Y,W)\mid Z$. No
positivity is required of $C$.

\textbf{Proof.} Fix values $z$, $y$, $w$. By positivity $P(y,w,z)>0$ for every
$y$ and $w$ in their value sets, so all conditionals below are defined. The
first hypothesis gives $P(c\mid y,z,w)=P(c\mid z,w)$ and the second gives
$P(c\mid y,z,w)=P(c\mid z,y)$. Hence $P(c\mid z,w)=P(c\mid z,y)$ for all
$y,w$. The left side does not depend on $y$ and the right side does not
depend on $w$, so both equal a function $q(c,z)$ only. Then
$P(c\mid z)=\sum_{y,w}P(c\mid y,z,w)P(y,w\mid z)=q(c,z)\sum_{y,w}P(y,w\mid z)
=q(c,z)$, and therefore $P(c\mid y,z,w)=q(c,z)=P(c\mid z)$ for all $y,w$, which
is $C\perp(Y,W)\mid Z$. \hfill$\square$

\textbf{Proof of Proposition~\ref{prop:unique}.} Let $S_1,S_2$ be blankets. Put
$Z=S_1\cap S_2$, $A=S_1\setminus S_2$, $B=S_2\setminus S_1$ and
$R=U\setminus(S_1\cup S_2)$; these are disjoint and cover $U$.
Since $S_1=Z\cup A$ is a blanket, $C\perp(B,R)\mid(Z,A)$.
Since $S_2=Z\cup B$ is a blanket, $C\perp(A,R)\mid(Z,B)$, and by weak union
$C\perp A\mid(Z,B,R)$. Apply the Lemma with $Y=A$ and $W=(B,R)$: the two
hypotheses are $C\perp A\mid(Z,B,R)$ and $C\perp(B,R)\mid(Z,A)$, both
established. Hence $C\perp(A,B,R)\mid Z$, i.e.\ $Z=S_1\cap S_2$ is a blanket.
Now let $S_1,S_2$ be inclusion-minimal blankets. Then $S_1\cap S_2$ is a
blanket contained in both, so by minimality $S_1=S_1\cap S_2=S_2$.
\hfill$\square$

\textbf{Link to Definition~\ref{def:parents}.} If $C$ is a function of $U$ (the exact
case of Condition 2, where some subset of $U$ is $0$-sufficient), then for
any $S\subseteq U$, $S$ is a blanket iff $\Hent(C\mid S)=0$. Indeed, if
$\Hent(C\mid S)=0$ then $C$ is a function of $S$ and is independent of
everything given $S$; conversely, if $C\perp U\setminus S\mid S$ then
$\Hent(C\mid S)=\Hent(C\mid S,U\setminus S)=\Hent(C\mid U)=0$. So under
positivity the exact parent set of Definition~\ref{def:parents} is unique.

\textbf{Proof of Proposition~\ref{prop:welldef}.} Nodes are equivalence classes of
accepted candidates, i.e.\ distinct partitions of $X$; since the accepted set
is fixed (assumption (c)), the node set is determined. Each class is assigned
the earliest site of its members under the fixed linear extension, which is
determined. For each node $C$ the candidate parents form a positive
collection by (b), so its exact parent set is unique by the result just
proved; hence the edge set is determined. Every edge goes from a node whose
site is strictly earlier to a node whose site is later in a fixed linear
order, so a directed cycle would require a site strictly earlier than
itself; the graph is acyclic. Finally, by Proposition~\ref{prop:reparam}, replacing a
value map by an invertible reparametrisation leaves the partition, hence the
node, and all conditional-independence statements, hence the edges,
unchanged. \hfill$\square$

\textbf{Caveat recorded with the proof.} Positivity fails whenever the
candidate set contains two members one of which is a function of the other
(for example a node together with its own parents), and equivalent
candidates violate it trivially; merging (Definition~\ref{def:equiv}) is therefore a
precondition, not a convenience. When positivity fails, distinct minimal
parent sets can exist (for $C=C_3=C_1\oplus C_2$ both $\{C_3\}$ and
$\{C_1,C_2\}$ are minimal). The programme then reports the set of all
minimal parent sets and every downstream quantity as a range over the
resulting graphs; the choice of a canonical candidate set (for example, the
latest-site antichain that is sufficient) is an open design decision for
Phase 1 and is tested on planted system V1.
